\section{De la variation au capital : composer correctement les rendements}

Si une unité monétaire est investie et réinvestie, sa valeur après $T$ périodes est
\[
W_T=W_0\prod_{t=1}^{T}(1+R_t).
\]

Pour $W_0=1$, la courbe $W_T$ est souvent appelée \emph{wealth curve} ou courbe de richesse. Elle ne constitue pas encore une simulation rétrospective complète de stratégie (backtest), mais elle apprend déjà à respecter la composition multiplicative des gains et pertes.

\coursefigure{figures/fig09_wealth_curve.pdf}{Courbe de richesse obtenue en composant les rendements simples.}{fig:wealth}

\begin{attentionbox}{Additionner des rendements simples sur de longues périodes}
La somme $\sum_t R_t$ peut fournir une approximation pour de petites variations, mais elle ne représente pas exactement la croissance composée. Pour reconstituer un capital, utiliser le produit des facteurs $(1+R_t)$ ou l'exponentielle de la somme des rendements logarithmiques.
\end{attentionbox}

\section{Qualité des données : le premier ``modèle'' est un contrôle}

Une chaîne de traitement (pipeline) de trading ne devrait jamais commencer par \texttt{model.fit}. Il devrait commencer par des contrôles de cohérence. Les erreurs les plus dangereuses sont celles qui ne provoquent aucune exception Python mais modifient silencieusement le sens économique des données.

\subsection{Contrôles minimaux}

Pour une table OHLCV, on vérifiera au minimum :
\begin{enumerate}
    \item dates triées dans l'ordre croissant ;
    \item absence de doublons non expliqués ;
    \item prix strictement positifs lorsque la convention l'exige ;
    \item $High\geq\max(Open,Close)$ ;
    \item $Low\leq\min(Open,Close)$ ;
    \item $Ask\geq Bid$ si des quotes sont présentes ;
    \item type et unité du volume ;
    \item taux de valeurs manquantes ;
    \item calendrier attendu et trous anormaux ;
    \item provenance et version de la source.
\end{enumerate}

\begin{lstlisting}
assert df["date"].is_monotonic_increasing
assert not df["date"].duplicated().any()
assert (df["high"] >= df[["open", "close"]].max(axis=1)).all()
assert (df["low"]  <= df[["open", "close"]].min(axis=1)).all()
\end{lstlisting}

Un \texttt{assert} ne remplace pas un audit de donnée, mais il transforme une hypothèse silencieuse en condition explicite.

\subsection{Valeur aberrante ou événement réel ?}

En analyse de données classique, une valeur très extrême, souvent qualifiée de valeur aberrante (\emph{outlier}), est parfois traitée comme une anomalie à supprimer. En finance, cette réflexion est dangereuse : un crash, un gap ou une annonce importante est précisément un événement économique que le modèle doit parfois apprendre à gérer.

Avant de supprimer un point extrême, il faut demander :
\begin{itemize}
    \item s'agit-il d'une erreur de collecte ou d'un mouvement réel ?
    \item la source secondaire confirme-t-elle l'observation ?
    \item une action de société explique-t-elle la rupture ?
    \item l'événement est-il hors des horaires habituels ?
\end{itemize}

\begin{engineerbox}{Tracer une décision de nettoyage}
Dans un projet sérieux, toute suppression ou correction de donnée doit être reproductible. Conserver le fichier brut en lecture seule, produire un fichier nettoyé par script, et journaliser les transformations importantes.
\end{engineerbox}

\section{Premier pont vers le Machine Learning}

Le Machine Learning (ML, apprentissage automatique) sera réactivé progressivement lorsque ses notions deviendront utiles. À ce stade, trois idées suffisent : une observation, des variables explicatives et une cible à prédire.

\begin{recallbox}{Feature, Target, Prédiction}
En apprentissage supervisé, une observation au temps $t$ est représentée par un vecteur de variables $\mathbf{x}_t$. Le modèle apprend une fonction
\[
\widehat{y}_t=f_\theta(\mathbf{x}_t).
\]
Les composantes de $\mathbf{x}_t$ sont les variables explicatives, ou \emph{features}. La variable $y_t$ est la cible, ou \emph{target}. En trading, l'ordre temporel impose une contrainte supplémentaire : chaque feature doit être disponible au moment de la décision.
\end{recallbox}

Supposons par exemple que l'on veuille prédire le signe du rendement de demain :
\[
y_t=\mathbf{1}(r_{t+1}>0).
\]
La cible utilise le futur par définition ; c'est normal, car elle sert à apprendre. Mais aucune feature de $\mathbf{x}_t$ ne doit dépendre de $r_{t+1}$.

On pourrait construire :
\[
\mathbf{x}_t=
\left[
 r_t,
 r_{t-1},
 \widehat{\sigma}_{t,20},
 \text{volume}_t,
 \text{spread}_t
\right].
\]
Toutes ces quantités sont observables au temps $t$, sous réserve d'une convention claire sur l'heure exacte de calcul.

\begin{attentionbox}{Le futur peut se cacher dans une transformation apparemment innocente}
Une moyenne glissante centrée, une mise à l'échelle (\emph{scaling}) calculée sur tout le jeu de données ou une colonne ``next return'' mal décalée peuvent faire entrer de l'information future dans les features. Le modèle semble alors excellent parce qu'il a accès à une information qui n'aurait pas été disponible au moment réel de la décision. Nous consacrerons un chapitre entier à ce problème.
\end{attentionbox}

\section{Premier pont vers le Reinforcement Learning}

